\cotefile{44}
\scanpage{2}\mspage{\unsure{1|7}}
\begin{center}\uline{Programme de Travail}\end{center}

\deleted{\unsure{B|3}}\uline{A}

% Right-hand square bracket drawn to the left of Th. 1.
\noindent\llap{\smash{\raisebox{-2.2ex}{$\left.\rule{0pt}{4.8ex}\right]$}}\quad}\uline{Th. 1}\qquad $N(\unsure{\zeta|\xi})\subset\operatorname{Alb}(S)$ n'est autre \inserted{\deleted{des}ensemblistement} que le noyau de l'application naturelle
\[
\varphi_{\zeta}:\quad\operatorname{Alb}(\unsure{s|S})\longrightarrow
\mathcal{Z}^{i}(X)/\mathcal{Z}^{i}_{\mathrm{alb}}(X)
\]

% Two overlapping right brackets in the left margin.
\noindent\llap{\smash{\raisebox{-1.5ex}{$\left.\rule{0pt}{3.6ex}\right]$}}\quad}\llap{\smash{\raisebox{-3.3ex}{$\left.\rule{0pt}{3.1ex}\right]$}}\quad}Th. 2\qquad La dimension de $\operatorname{Alb}(S)/N(\zeta)=P_{\zeta}$ est majorée par un nb ne dépendant que de $X$ [en fait, par $\frac{1}{2}b^{2i-1}(X)$ !]

Ces deux th. impliquent

\noindent\llap{\smash{\raisebox{-2.2ex}{$\left.\rule{0pt}{4.8ex}\right]$}}\quad}\uline{Th. 3}\quad $Z^{i}_{\mathrm{alg}}(X)/\mathcal{Z}^{i}_{\tikz[baseline=(a.base)]{\node[ellipse,draw,inner sep=0.3pt] (a) {$\scriptstyle\mathrm{alb}$};}}(X)=P^{i}(X)$ a une structure naturelle de V.A., $\deleted{\unsure{P}}$et $P^{i}(X)$ et $P^{\unsure{2m-i|n+1-i|n-i}}(X)$ sont en dualité canonique

% slanted; circled note with an arrow to Th. 4
\margin{\tikz[baseline=(n.base)]{\node[draw,ellipse,inner sep=1pt,align=center] (n) {\unsure{Canulé|Annulé|Cancelé}\\tel quel};\draw[->] (n.south) -- ++(0.18,-0.2);}}
% slanted
\margin{\textquotedblleft{}\uline{Hodge}}
\uline{Th. 4}\quad Si $n=2m-1$, alors l'autodualité canonique de $P^m(X)$ [donc \unsure{la|sa} corr. sur $\underline{P}^m(X)$] est \unsure{associé|associée} à une \uline{polarisation} de $P^m(X)$
% archive sheet number: 1
